\documentclass[journal]{IEEEtran}

\usepackage{amsmath}
\usepackage{booktabs}
\usepackage{graphicx}
\usepackage{url}

\title{Leakage-Controlled Evaluation of Feature Stability and Selective Reliability in CARDIODIAG}

\author{[AUTHOR NAMES AND AFFILIATIONS TO BE ADDED]}

\begin{document}
\maketitle

\begin{abstract}
This manuscript reports a leakage-controlled reproduction and extension of the
CARDIODIAG cardiovascular disease prediction framework. The study separates the
historical notebook protocol from a corrected pipeline in which imputation,
encoding, outlier handling, feature selection, model selection, and uncertainty
calibration respect data boundaries. A stability-aware training-only feature
selection component is evaluated against corrected logistic, voting, and
stacking baselines, and a separate selective-reliability layer measures whether
confidence ranks correct decisions above errors. Internal nested-validation
results are generated directly from machine-readable experiment artifacts. No
independent external validation was available, so all reported evidence is
internal development/replication evidence and the absence of external
validation is treated as a limitation.
\end{abstract}

\begin{IEEEkeywords}
cardiovascular disease, leakage control, nested validation, explainable AI,
calibration, conformal prediction, selective prediction.
\end{IEEEkeywords}

\section{Introduction}
Cardiovascular prediction models must be evaluated for discrimination,
calibration, robustness, and transportability rather than accuracy alone.
CARDIODIAG provides an important starting point because it combines explainable
feature reduction with optimized ensemble classifiers \cite{khan2026cardiodiag}.
This study asks two falsifiable questions under properly separated validation:
whether a resampled feature selector is more stable than the same relevance
score applied once, and whether an abstention layer can identify higher-risk
predictions. The latter is treated as empirical selective reliability rather
than evidence that a model semantically understands uncertainty.

\textbf{Contribution status.} The proposed contribution is a
leakage-controlled stability--reliability decoupling audit for CARDIODIAG, not
a confirmed novel algorithm. It tests whether more stable selected feature sets
actually yield better discrimination, calibration, or selective reliability.
The current evidence is internal and does not establish clinical utility,
deployment readiness, or external validity.

\section{Related Work and Research Gap}
Feature-selection stability is an established property of a selection procedure,
not a substitute for predictive validation. Nogueira \emph{et al.} provide a
chance-corrected stability statistic for comparing selected sets
\cite{nogueira2018stability}, while resampling-based stability selection has
distinct assumptions and error-control goals \cite{meinshausen2010stability}.
Nested validation is required when the same data support model selection and
performance estimation \cite{varma2006bias,cawley2010overfitting}. Calibration,
discrimination, and selective risk are likewise distinct properties of a
prediction model \cite{vancalster2019calibration}. Conformal prediction can
provide marginal coverage under exchangeability, but coverage should not be
assumed under population shift \cite{angelopoulos2023conformal}.

The gap is narrower than a new feature-selection algorithm: this study tests
whether training-only feature-set stability in the CARDIODIAG representation is
associated with held-out predictive and selective reliability. The reporting
and risk-of-bias framing follows TRIPOD+AI and PROBAST+AI
\cite{collins2024tripodai,moons2025probastai}.

\section{Audit of the Original Implementation}
The supplied repository is notebook-centric and contains multiple historical
protocols. The historical reproduction and its deviations are documented in
\texttt{docs/baseline\_reproduction\_report.md}. In particular, available
notebooks contain paths where resampling, preprocessing, and dimensionality
reduction occur before the evaluation split.

\subsection{Historical Reproduction Claims}
The original paper's reported protocol cannot be reproduced bit-for-bit from
the supplied notebooks: split sizes, preprocessing order, model blocks, search
grids, and saved tables are inconsistent. The native reproduction path
specified in the audit applied SMOTE before an 80/20 split without
stratification; on \texttt{heart5.csv} it produced 1,016 rows after SMOTE and
reported OSEC accuracy 0.8873, ROC-AUC 0.9516, and Brier score 0.0890, while
OVEC reported accuracy 0.8627, ROC-AUC 0.9163, and Brier score 0.1236. These
are protocol-specific reproduction artifacts, not leakage-controlled
generalization estimates. The corrected single-split comparison reported
OSEC accuracy 0.8957, ROC-AUC 0.9348, and Brier score 0.0875, with stable OSEC
accuracy 0.9000, ROC-AUC 0.9360, and Brier score 0.0864; these values are
not the primary corrected evidence because they use one split. Accordingly,
the original paper's numerical claims are not treated as validated
corrected-pipeline results.

\section{Methods}
\subsection{Corrected Pipeline}
All data-dependent operations are fitted on the training data of the relevant
fold. Feature types are declared in the local dataset manifest: nominal and
ordinal coded predictors are one-hot encoded, identifiers are excluded, and
numeric missingness/outlier treatment is training-only. Unknown external data
must provide an explicit feature mapping before the external gate can pass.

\subsection{Stability-Aware Feature Reduction}
For feature $j$ and inner resample $b$, let $r_{j,b}$ denote the combined
training-only relevance score and $s_{j,b}$ its selection indicator. Relevance
is estimated after splitting each resample into a fitting part and a validation
part; permutation importance is evaluated on the latter, which remains inside
the outer training boundary. We define
\begin{equation}
\pi_j = \frac{1}{B}\sum_{b=1}^{B}s_{j,b}, \qquad
\bar r_j = \frac{1}{B}\sum_{b=1}^{B}r_{j,b}.
\end{equation}
Features are retained when their selection frequency and mean relevance pass
the prespecified thresholds. To test whether this improves stability, the
selected sets are compared with a single-pass relevance selector over matched
repeated stratified training perturbations using the chance-corrected stability
statistic of Nogueira \emph{et al.} \cite{nogueira2018stability}. The reported
interval is a paired bootstrap interval over repeat-level stability estimates.
This is an operational stability rule, not a causal interpretation of feature
importance.

\subsection{Uncertainty and Calibration}
The project reports Brier score, reliability-bin error, calibration intercept
and slope, split-conformal prediction-set coverage, abstention rate, confidence-
error AUC, area under the risk-coverage curve (AURC), and selective risk. The
confidence-error AUC tests whether confidence ranks correct predictions above
errors; it is not a probability of correctness. Conformal coverage is
interpreted under its exchangeability assumption \cite{angelopoulos2023conformal}
and is measured again under any available dataset shift.
The implementation uses ten uniform bins for ECE, omits empty bins, and
averages the remaining absolute calibration gaps equally. Its AURC value is
the arithmetic mean of cumulative error risks over retained-prefix sizes,
rather than a numerically integrated curve; this implementation detail limits
direct comparison with studies using a different AURC convention.

\section{Experimental Protocol}
The historical protocol, corrected common-split protocol, and nested-validation
protocol are kept separate. Hyperparameter selection is performed only in inner
folds. The final internal holdout is not repeatedly consulted. Reporting follows
the transparency principles of TRIPOD+AI \cite{collins2024tripodai} and the
development/evaluation distinctions in PROBAST+AI \cite{moons2025probastai}.

For feature-selection stability, five stratified folds were repeated three
times. The stable selector was compared with a one-shot selector using the
recorded training-fold feature selections. The standalone corrected stability artifacts used eight relevance resamples,
whereas the stable selector embedded in the nested and linkage pipelines used
six. For selective reliability and the paired stability--reliability analysis,
ten random seeds (42--51) used a 56.25\%
fitting, 18.75\% calibration, and 25\% evaluation split. Evaluation labels were
not used for fitting or calibration. These bounded repeated analyses are
exploratory and are not substitutes for a full confirmatory evaluation.

The primary nested endpoint is outer-fold ROC-AUC for the listed candidates;
the stability--reliability linkage is an exploratory paired comparison of
stable OSEC against corrected OSEC. PR-AUC, Brier score, calibration,
feature stability, and selective metrics are secondary. The reported intervals
are descriptive bootstrap summaries over the recorded folds, repeat-level
stability summaries, or seeds; they are not independent-patient confidence
intervals.

As a separate exploratory benchmark, the 70,000-row Kaggle \texttt{cardio}
task was evaluated with three shared stratified outer folds, two inner folds,
seed 42, and bounded inner searches for Logistic regression, Random Forest,
HistGradientBoosting, XGBoost, LightGBM, and CatBoost. This task has a
different target and schema and is not pooled with the two CARDIODIAG
benchmarks. A fixed one-third out-of-fold ensemble of XGBoost,
HistGradientBoosting, and Logistic regression was evaluated on the same outer
folds.

\section{Datasets and Provenance}
The supplied Five States and Mendeley datasets are development and replication
benchmarks because they were already used in the base study. The Mendeley
dataset is attributed to Doppala and Bhattacharyya~\cite{doppala2021dataset}
(DOI: \url{https://doi.org/10.17632/dzz48mvjht.1}). Dataset hashes and their
assigned roles are recorded in
\texttt{data/manifests/local\_datasets.json}. Any dataset intended for
independent external validation must be identified and frozen in
\texttt{data/manifests/external\_dataset.json} before label-based modeling
decisions are made. No independent external validation dataset was available
for the present evaluation.

\section{Results}
Table~\ref{tab:nested} reports the post-boundary-corrected nested-validation
artifacts. It is internal evidence only and must not be described as
independent external evidence. Earlier historical and schema-corrected
artifacts remain available under their original prefixes and are not silently
replaced.

\begin{table*}[t]
\centering
\caption{Nested internal validation summary.}
\label{tab:nested}
\begin{tabular}{llrrrrr}
\toprule
Dataset & Candidate & ROC-AUC & PR-AUC & Brier & ECE & ROC-AUC CI \\
\midrule
heart5.csv & corrected\_logistic & 0.928 $\pm$ 0.006 & 0.935 & 0.107 & 0.104 & [0.907, 0.942] \\
heart5.csv & corrected\_ovec & 0.930 $\pm$ 0.007 & 0.934 & 0.100 & 0.109 & [0.909, 0.944] \\
heart5.csv & corrected\_osec & 0.926 $\pm$ 0.010 & 0.929 & 0.109 & 0.156 & [0.899, 0.936] \\
heart5.csv & stable\_logistic & 0.928 $\pm$ 0.006 & 0.935 & 0.107 & 0.101 & [0.907, 0.942] \\
heart5.csv & stable\_osec & 0.925 $\pm$ 0.011 & 0.929 & 0.109 & 0.156 & [0.896, 0.934] \\
MendeleyCardiovascular\_Disease\_Dataset.csv & corrected\_logistic & 0.994 $\pm$ 0.003 & 0.995 & 0.028 & 0.172 & [0.988, 0.996] \\
MendeleyCardiovascular\_Disease\_Dataset.csv & corrected\_ovec & 0.996 $\pm$ 0.001 & 0.997 & 0.026 & 0.136 & [0.993, 0.998] \\
MendeleyCardiovascular\_Disease\_Dataset.csv & corrected\_osec & 0.995 $\pm$ 0.002 & 0.996 & 0.024 & 0.211 & [0.992, 0.998] \\
MendeleyCardiovascular\_Disease\_Dataset.csv & stable\_logistic & 0.994 $\pm$ 0.003 & 0.995 & 0.028 & 0.188 & [0.989, 0.996] \\
MendeleyCardiovascular\_Disease\_Dataset.csv & stable\_osec & 0.996 $\pm$ 0.002 & 0.997 & 0.024 & 0.216 & [0.990, 0.997] \\
\bottomrule
\end{tabular}

\end{table*}

\subsection{Nested Validation}
The corrected nested runs used three outer folds, one outer repeat, and
three-fold inner selection with ROC-AUC scoring. Corrected OVEC had the
highest listed fold-mean ROC-AUC on both datasets: 0.9295 on Five States and
0.9957 on Mendeley. Stable OSEC was 0.9253 versus corrected OSEC 0.9259 on
Five States, and 0.9942 versus 0.9943 on Mendeley. The small differences do
not establish a uniform predictive advantage. Stable-candidate threshold
tuning occurred inside inner folds over 0.5, 0.75, and 1.0; these results
therefore are not numerically interchangeable with the fixed-threshold
stability experiment.

\subsection{Reliability}
The linkage artifacts compare stable OSEC with corrected OSEC over ten seeds.
Confidence-error AUC was above 0.5 for both candidates on both benchmarks, but
the paired differences did not consistently favor the stable selector.
Conformal coverage was measured on the evaluation partition and is not a
guarantee under shift. The linkage intervals and rank correlations are
descriptive and exploratory.

\begin{table*}[t]
\centering
\caption{Ten-seed internal selective-reliability summaries from the
post-boundary-corrected linkage artifacts. Paired differences and descriptive
bootstrap intervals are reported in Table~\ref{tab:linkage}; these are not
independent external-validation estimates.}
\label{tab:uncertainty}
\begin{tabular}{llrrrr}
\toprule
Dataset & Candidate & Conf.-error AUC & AURC & Conformal coverage & ECE \\
\midrule
Five States & Corrected OSEC & 0.762 & 0.0505 & 0.914 & 0.122 \\
Five States & Stable OSEC & 0.760 & 0.0516 & 0.915 & 0.121 \\
Mendeley & Corrected OSEC & 0.905 & 0.0041 & 0.914 & 0.221 \\
Mendeley & Stable OSEC & 0.905 & 0.0042 & 0.914 & 0.217 \\
\bottomrule
\end{tabular}

\end{table*}

\subsection{Feature-Selection Stability}
The post-boundary-corrected stability comparison is shown in
Table~\ref{tab:stability}. The stable-minus-one-shot differences were 0.0974
(95\% descriptive bootstrap interval [-0.0272, 0.1978]) on Five States and
0.0687 ([0.0000, 0.1671]) on Mendeley. The Mendeley interval touches zero,
and the artifact's prespecified positive-evidence flag is false. These
intervals are based on only three repeat-level differences and should be
interpreted as limited comparative evidence, not patient-level confidence
intervals or proof of generality.

\begin{table}[t]
\centering
\caption{Post-boundary-corrected feature-selection stability. $\Phi$ is the
Nogueira chance-corrected stability statistic; intervals summarize the three
repeat-level paired differences and are exploratory.}
\label{tab:stability}
\begin{tabular}{lrrrr}
\toprule
Dataset & Stable $\Phi$ & One-shot $\Phi$ & Difference & 95\% bootstrap CI \\
\midrule
Five States & 0.112 & 0.014 & 0.097 & [-0.027, 0.198] \\
Mendeley & 0.284 & 0.215 & 0.069 & [0.000, 0.167] \\
\bottomrule
\end{tabular}

\end{table}

\subsection{Stability--Reliability Linkage}
Table~\ref{tab:linkage} tests whether a more stable selected set translated into
better held-out reliability on the same repeated splits. Differences are
stable OSEC minus corrected OSEC; positive AURC or ECE differences indicate
worse selective risk or calibration, respectively. The paired intervals include
zero for all listed comparisons. On Mendeley, ROC-AUC, AURC, and
full-coverage risk favored corrected OSEC, while the other intervals included
zero. Thus, selector stability did not act as a consistent surrogate for
predictive or selective reliability.

\begin{table*}[t]
\centering
\caption{Paired stability--reliability linkage. $\Phi$ is cross-seed selected-
set stability. $\Delta$ values are stable OSEC minus corrected OSEC means; the
full bootstrap intervals are stored in the machine-readable linkage artifacts.}
\label{tab:linkage}
\begin{tabular}{lrrrrr}
\toprule
Dataset & Cross-seed $\Phi$ & $\Delta$ ROC-AUC & $\Delta$ Conf.-AUC & $\Delta$ AURC & $\Delta$ ECE \\
\midrule
Five States & 0.201 & -0.0009 & -0.0015 & +0.0011 & -0.0001 \\
Mendeley & 0.178 & -0.0002 & -0.0001 & +0.0002 & -0.0045 \\
\bottomrule
\end{tabular}

\end{table*}

\section{Discussion}
The nested results do not support a uniform performance advantage for the
stability-aware selector. On the Five States benchmark, corrected OVEC has the
highest mean ROC-AUC among the listed candidates, while stable OSEC is lower.
On the Mendeley benchmark, corrected OVEC and stable OSEC are similar. The
post-boundary-corrected stability comparison does not establish a positive
stability advantage: its Five States interval crosses zero and its Mendeley
interval touches zero. The ten-seed reliability results do not show
consistent improvement across discrimination, selective risk, coverage, and
calibration. These findings separate feature-set stability from predictive
reliability: more repeatable feature selection does not by itself establish
better calibrated or safer predictions. A selector is not justified merely
because it reduces dimensionality or includes an explainability component.

\section{External Validation and Threats to Validity}
No independent compatible dataset was available for a defensible frozen
external evaluation within this project. The supplied Mendeley dataset is not
new external evidence because it was used in the base paper, and the public UCI
component sources overlap the Five States composite. Consequently, this study
reports internal nested validation and reliability analysis only, and does not
claim independent external validation. This caution is consistent with the
documented heterogeneity and performance degradation commonly observed when
cardiovascular prediction models are evaluated outside their development data
\cite{wessler2021external}.

\section{Statistical Analysis and Robustness}
The study prespecifies ROC-AUC as the primary comparison and reports PR-AUC,
classification metrics, Brier score, calibration intercept and slope, uniform
ten-bin ECE, selected-set stability, confidence-error AUC, AURC, and conformal
coverage as secondary quantities. Paired comparisons use identical partitions
for the candidates. No multiplicity-adjusted confirmatory claim is made across
the secondary metrics. Synthetic missingness or measurement perturbations,
when run, will be reported as stress tests and not as independent validation.

The post-boundary-corrected stability runs found differences of 0.0974
[-0.0272, 0.1978] on Five States and 0.0687 [0.0000, 0.1671] on Mendeley.
The ten-seed linkage runs found no consistent reliability advantage: the paired
ROC-AUC difference was -0.0009 [-0.0027, 0.0007] on Five States and -0.0002
[-0.0004, -0.0001] on Mendeley, while calibration and selective metrics did
not uniformly favor the stable candidate. These are bounded internal results
and remain subject to dependence, small effective replication counts, and
multiple comparisons.

The separate Kaggle benchmark produced pooled outer ROC-AUC values of 0.80043
for CatBoost and 0.80022 for XGBoost in the tuned run. The prior XGBoost
benchmark was 0.80038. The fixed uniform ensemble achieved 0.79897 and
therefore did not improve the standalone boosted models. The CatBoost
difference is marginal, exploratory, and not evidence of practical or
statistical superiority.

A compact candidate-pool rerun evaluated Elastic Net, RBF SVM, KNN, Extra
Trees, Gaussian NB, and LDA alongside the scalable boosting candidates. On
70,000 rows, RBF SVM and KNN were prespecified as computationally infeasible
and were not run. In that separate protocol, stacking achieved pooled
ROC-AUC 0.80028, XGBoost 0.80008, CatBoost 0.79998, and fixed soft voting
0.79943. The small differences do not establish an ensemble or model
superiority.

An evaluation-only stress test at seed 42 applied 10\% independent missingness
and numeric noise with standard deviation equal to 5\% of each training feature
standard deviation. On Mendeley, ROC-AUC fell from 0.998 to 0.810 under the
numeric-noise perturbation for corrected OSEC and from 0.998 to 0.815 for stable
OSEC. On Five States, 10\% missingness reduced ROC-AUC from approximately 0.935
to 0.922 for both candidates. These are synthetic single-split stress tests,
not independent external validation.

\section{Limitations and Ethical Considerations}
The study is retrospective, benchmark-based, and not a clinical diagnostic
device. It does not establish prospective utility, patient benefit, fairness
across unsupported demographic groups, or clinical safety. Dataset provenance,
license conditions, outcome definitions, and possible source overlap constrain
transportability. The full five-candidate custom nested protocol was not
completed on the 70,000-row Kaggle task because of computational cost; a
reduced corrected-logistic versus stable-logistic analysis is labeled
separately.

\section{Reproducibility}
The code, tests, dependency lock, data manifests, protocol reports, and result
generators are maintained in the repository. Raw restricted data must not be
committed.

\section{Conclusion}
This work delivers a reproducible leakage-controlled audit and a falsifiable
stability--reliability decoupling audit for CARDIODIAG. Internal evidence shows
that feature-set stability and selective reliability can move independently:
the corrected stability comparison does not establish a stable-selector
advantage, and the stable selector does not produce a consistent predictive or
reliability benefit.
Independent external and prospective evaluation remain necessary before any
generalization or clinical-readiness claim.

\bibliographystyle{IEEEtran}
\bibliography{references}
\end{document}